% ══════════════════════════════════════════════════════════════════════════════
%  Chapter 17: Desktop Admin Interface (Tauri)
% ══════════════════════════════════════════════════════════════════════════════
\chapter{Desktop Admin Interface (Tauri)}
\label{ch:desktop-tauri}

\begin{learnbox}
\begin{itemize}
  \item Build the same admin functionality using Tauri 2 with a Rust backend and React frontend
  \item Implement IPC commands that bridge Rust and TypeScript across the process boundary
  \item Compare the Tauri architecture with the Iced MVU approach from the previous chapter
  \item Use Zustand for state management and React Query for data fetching in the frontend
\end{itemize}
\end{learnbox}

\lettrine{T}{auri} is a \index{framework}framework for building \index{desktop application}desktop applications that pair a \index{Rust!backend}Rust backend with a \index{web-technology}web-technology \index{frontend}frontend (HTML/CSS/JS). Unlike \index{Electron}Electron, \index{Tauri}Tauri does not bundle Chromium---it uses the operating system's native \index{webview}webview, resulting in much smaller \index{binary}binaries (typically 5--24 MB vs. 150+ MB). The \index{Rust!backend}Rust side exposes \emph{\index{command}commands} via \code{\#[\index{tauri::command}tauri::command]}, which the \index{frontend}frontend calls through an \code{\index{invoke}invoke()} function over an \index{IPC}IPC bridge. The result is a two-language \index{architecture}architecture: \index{Rust!backend}Rust handles business \index{logic}logic and \index{system}system access, while \index{React}React handles the \index{UI}UI.

\begin{notebox}[Architectural Comparison]
This chapter builds the same admin interface as Chapter~\ref{ch:desktop-iced}, but with Tauri 2---a Rust backend connected to a React/TypeScript frontend via IPC. If you read Chapter~\ref{ch:desktop-iced} first, you will see exactly where the architectures diverge: state management, rendering, and the boundary between business logic and UI.
\end{notebox}

This chapter explains the \code{\index{bitcoin-desktop-ui-tauri}bitcoin-desktop-ui-tauri} \index{Cargo!crate}crate as a \index{Rust!developer}Rust developer reads it: \textbf{how a \index{Tauri}Tauri 2 app bridges a \index{Rust!backend}Rust backend to a \index{React}React/\index{TypeScript}TypeScript \index{frontend}frontend}, \textbf{how \index{command}commands flow across the \index{IPC}IPC boundary}, and \textbf{how the \index{frontend}frontend calls our \index{Bitcoin}Bitcoin \index{implementation}implementation's \index{admin dashboard}admin API} through that bridge.

\section{What This UI Is (In One Sentence)}
\label{sec:ch17-oneline}

The \index{Tauri}Tauri \index{desktop application}desktop \index{admin dashboard}admin UI is a \textbf{two-language \index{desktop application}desktop \index{application}application}: a \index{Rust!backend}Rust \index{backend}backend exposes 22 \code{\#[\index{tauri::command}tauri::command]} handlers that wrap our \code{\index{bitcoin-api}bitcoin-api} \index{Cargo!crate}crate, while a \index{React}React/\index{TypeScript}TypeScript \index{frontend}frontend calls those handlers through \index{Tauri}Tauri's \index{IPC}IPC bridge using \code{\index{invoke}invoke()}.

The Tauri desktop admin UI consists of: \code{main} (Tauri app setup, state management, command registration), \code{BitcoinApiService::*} (service layer wrapping \code{AdminClient} HTTP calls), \code{commands::*} (22 Tauri command handlers across 6 modules), \code{commands.ts} \(\to\) \code{invoke()} (TypeScript-to-Rust bridge), and React pages (18 views), components (14 shared), hooks (2 custom).

% ──────────────────────────────────────────────────────────────────────────────
\section{Why Tauri? (Iced vs. Tauri --- An Architectural Comparison)}
\label{sec:ch17-comparison}

In Chapter~\ref{ch:desktop-iced} we built the same \index{admin dashboard}admin \index{dashboard}dashboard with \index{Iced}Iced---a pure-\index{Rust!backend}Rust, single-language \index{MVU}MVU \index{framework}framework. Chapter~\ref{ch:desktop-tauri} builds the identical \textbf{\index{feature}feature set} (\index{blockchain}blockchain info, \index{wallet}wallet \index{management}management, \index{mining}mining controls, \index{health check}health checks) with \index{Tauri}Tauri, which gives us a \textbf{\index{Rust!backend}Rust \index{backend}backend paired with a \index{web}web \index{frontend}frontend}. Both approaches have real \index{trade-off}trade-offs:

\begin{table}[htbp]
\caption{Iced vs. Tauri: Architectural comparison}
\label{tab:iced-vs-tauri}
\centering\small
\begin{tabularx}{0.95\textwidth}{lXX}
\toprule
\textbf{Aspect} & \textbf{Iced (Chapter~\ref{ch:desktop-iced})} & \textbf{Tauri (this chapter)} \\
\midrule
Languages & Rust only & Rust backend \texttt{+} TypeScript/React frontend \\
UI paradigm & MVU (Model-View-Update) & Component-based (React) with IPC bridge \\
State management & Single \code{AdminApp} struct & Zustand (client) \texttt{+} React Query (server) \\
Async boundary & Manual Tokio runtime in background thread (\code{spawn\_on\_tokio}) & Tauri handles async natively---\code{\#[tauri::command]} functions are async \\
Event vocabulary & \code{Message} enum with \(\sim\)50 variants & React Router routes \texttt{+} React Query cache keys \\
Styling & Iced built-in widgets (programmatic) & Tailwind CSS (declarative) \\
Shared code & \code{bitcoin-api} crate (Rust) & \code{bitcoin-api} crate (Rust)---same crate, same API calls \\
Binary size & Smaller (pure Rust) & Larger (embeds Chromium webview) \\
Ecosystem & Iced widgets only & Full npm ecosystem (React Query, Zod, Lucide, etc.) \\
\bottomrule
\end{tabularx}
\end{table}

The key point: both apps call exactly the same \code{AdminClient} methods from the \code{bitcoin-api} crate. The difference is where the UI layer lives (Rust vs. browser) and how events cross the boundary (Iced \code{Message} vs. Tauri \code{invoke}).

\begin{infobox}[Tip]
Both the Iced admin (Chapter~\ref{ch:desktop-iced}) and this Tauri admin consume the same \code{bitcoin-api} crate for blockchain communication. This means the backend API is framework-agnostic---you could add a third frontend without changing any backend code.
\end{infobox}

% ──────────────────────────────────────────────────────────────────────────────
\section{Diagram: Tauri IPC Architecture}
\label{sec:ch17-ipc-arch}

Data flows from the React frontend through IPC to the Rust backend and then to the blockchain API server:

\begin{minted}[fontsize=\footnotesize]{text}
+-------------------------------------------------------------------------+
|  TAURI APPLICATION                                                      |
|                                                                         |
|  +--------------------------+    IPC     +----------------------------+ |
|  |   React Frontend         |  invoke()  |   Rust Backend             | |
|  |   (WebView / Chromium)   | -------->  |   (src-tauri/)             | |
|  |                          |            |                            | |
|  |  Page Component          |            |  #[tauri::command]         | |
|  |    +- useQuery()         |            |    +- reads State<RwLock>  | |
|  |       +- commands.ts     |            |       +- BitcoinApiService | |
|  |          +- invoke(name) |  <-------- |          +- AdminClient    | |
|  |             {payload}    |  Result<V> |             +- HTTP call   | |
|  +--------------------------+            +-------------+--------------+ |
|                                                        |                |
+--------------------------------------------------------|----------------+
                                                         | HTTP
                                                         v
                                              +---------------------+
                                              |  Bitcoin API Server  |
                                              |  (Chapter 24)        |
                                              |  :8080               |
                                              +---------------------+
\end{minted}

% ──────────────────────────────────────────────────────────────────────────────
\section{File Organization}
\label{sec:ch17-files}

The Tauri application is split between two directories:

\begin{itemize}
  \item \textbf{Rust backend} (\code{src-tauri/}):
  \begin{itemize}
    \item \code{src/main.rs} (app setup, state, command registration)
    \item \code{src/api/service.rs} (service layer wrapping \code{AdminClient})
    \item \code{src/commands/} (22 handlers across 6 modules)
  \end{itemize}
  \item \textbf{React frontend} (\code{src/}):
  \begin{itemize}
    \item \code{src/commands.ts} (IPC bridge via \code{invoke()})
    \item \code{src/pages/} (18 views)
    \item \code{src/components/} (14 shared components)
    \item \code{src/hooks/} (custom hooks for queries)
  \end{itemize}
\end{itemize}

% ──────────────────────────────────────────────────────────────────────────────
\section{Rust Backend Architecture}
\label{sec:ch17-rust-backend}

The Rust backend follows a layered architecture:

\begin{enumerate}
  \item \textbf{Tauri state} (\code{main.rs}): Global state managed via \code{State<RwLock<AppState>>}.
  \item \textbf{Service layer} (\code{api/service.rs}): \code{BitcoinApiService} wraps \code{AdminClient} HTTP calls with async/await.
  \item \textbf{Command handlers} (\code{commands/*.rs}): 22 \code{\#[tauri::command]} functions that call the service and return results over IPC.
\end{enumerate}

\begin{listing}[htbp]
\caption{Main Tauri app setup (\code{src-tauri/src/main.rs})}
\label{lst:ch17-main}
\begin{minted}[fontsize=\footnotesize]{rust}
use tauri::State;
use std::sync::RwLock;

#[derive(Clone, serde::Serialize, serde::Deserialize)]
pub struct AppState {
    pub base_url: String,
    pub api_key: String,
}

fn main() {
    tauri::Builder::default()
        // Initialize AppState with default values
        .manage(AppState {
            base_url: "http://127.0.0.1:8080".to_string(),
            api_key: std::env::var("BITCOIN_API_ADMIN_KEY")
                .unwrap_or_else(|_| "admin-secret".to_string()),
        })
        // Register all command handlers
        .invoke_handler(tauri::generate_handler![
            cmd_get_blockchain_info,
            cmd_get_latest_blocks,
            cmd_get_blocks,
            cmd_get_block_by_hash,
            cmd_get_mining_info,
            cmd_generate_to_address,
            cmd_get_health,
            cmd_get_liveness,
            cmd_get_readiness,
            cmd_get_mempool,
            cmd_get_mempool_transaction,
            cmd_get_transactions,
            cmd_get_address_transactions,
            cmd_create_wallet_admin,
            cmd_get_addresses_admin,
            cmd_get_wallet_info_admin,
            cmd_get_balance_admin,
            cmd_send_transaction_admin,
            cmd_set_base_url,
            cmd_set_api_key,
            cmd_get_config,
        ])
        .run(tauri::generate_context!())
        .expect("error while running tauri application");
}
\end{minted}
\end{listing}

The state is accessed in command handlers via the \code{State<AppState>} parameter injected by Tauri.

\begin{listing}[htbp]
\caption{Service layer example (\code{src-tauri/src/api/service.rs})}
\label{lst:ch17-service}
\begin{minted}[fontsize=\footnotesize]{rust}
use bitcoin_api::{AdminClient, ApiConfig, ApiResponse, BlockchainInfo};
use serde_json::Value;

pub struct BitcoinApiService;

impl BitcoinApiService {
    pub async fn get_blockchain_info(
        base_url: String,
        api_key: String,
    ) -> Result<ApiResponse<BlockchainInfo>, String> {
        let config = ApiConfig {
            base_url,
            api_key: Some(api_key),
        };
        let client = AdminClient::new(config)
            .map_err(|e| e.to_string())?;
        client
            .get_blockchain_info()
            .await
            .map_err(|e| e.to_string())
    }

    pub async fn generate_to_address(
        base_url: String,
        api_key: String,
        address: String,
        nblocks: u32,
        maxtries: Option<u32>,
    ) -> Result<ApiResponse<Value>, String> {
        let config = ApiConfig {
            base_url,
            api_key: Some(api_key),
        };
        let client = AdminClient::new(config)
            .map_err(|e| e.to_string())?;
        client
            .generate_to_address(&address, nblocks, maxtries)
            .await
            .map_err(|e| e.to_string())
    }
}
\end{minted}
\end{listing}

\begin{listing}[htbp]
\caption{Command handler example (\code{src-tauri/src/commands/blockchain.rs})}
\label{lst:ch17-command}
\begin{minted}[fontsize=\footnotesize]{rust}
use crate::api::service::BitcoinApiService;
use tauri::State;
use serde_json::Value;

use crate::AppState;

#[tauri::command]
pub async fn cmd_get_blockchain_info(
    state: State<'_, AppState>,
) -> Result<Value, String> {
    let response = BitcoinApiService::get_blockchain_info(
        state.base_url.clone(),
        state.api_key.clone(),
    )
    .await?;

    // Return the entire API response as JSON
    serde_json::to_value(&response)
        .map_err(|e| e.to_string())
}

#[tauri::command]
pub async fn cmd_generate_to_address(
    state: State<'_, AppState>,
    address: String,
    nblocks: u32,
    maxtries: Option<u32>,
) -> Result<Value, String> {
    let response = BitcoinApiService::generate_to_address(
        state.base_url.clone(),
        state.api_key.clone(),
        address,
        nblocks,
        maxtries,
    )
    .await?;

    serde_json::to_value(&response)
        .map_err(|e| e.to_string())
}
\end{minted}
\end{listing}

\textbf{Walkthrough: the command handler pattern:}

Every command follows the same structure:

\begin{enumerate}
  \item Receive parameters from the IPC invocation.
  \item Call a corresponding \code{BitcoinApiService} method (which calls the bitcoin-api HTTP API).
  \item Return the result as JSON over IPC.
\end{enumerate}

This keeps the Rust side simple: it's a thin bridge between TypeScript and \code{AdminClient}.

% ──────────────────────────────────────────────────────────────────────────────
\section{TypeScript-to-Rust Bridge}
\label{sec:ch17-commands-ts}

\begin{listing}[H]
\caption{IPC command wrappers (\code{src/commands.ts})}
\label{lst:ch17-commands-ts}
\begin{minted}[fontsize=\footnotesize]{typescript}
import { invoke } from '@tauri-apps/api/core';
import { BlockchainInfo, BlockSummary, ApiResponse } from './types';

export async function getBlockchainInfo(): Promise<ApiResponse<BlockchainInfo>> {
    return invoke<ApiResponse<BlockchainInfo>>('cmd_get_blockchain_info');
}

export async function getLatestBlocks(): Promise<ApiResponse<BlockSummary[]>> {
    return invoke<ApiResponse<BlockSummary[]>>('cmd_get_latest_blocks');
}

export async function getBlocks(): Promise<ApiResponse<any>> {
    return invoke<ApiResponse<any>>('cmd_get_blocks');
}

export async function getBlockByHash(hash: string): Promise<ApiResponse<any>> {
    return invoke<ApiResponse<any>>('cmd_get_block_by_hash', { hash });
}

export async function getMiningInfo(): Promise<ApiResponse<any>> {
    return invoke<ApiResponse<any>>('cmd_get_mining_info');
}

export async function generateToAddress(
    address: string,
    nblocks: number,
    maxtries?: number,
): Promise<ApiResponse<any>> {
    return invoke<ApiResponse<any>>('cmd_generate_to_address', {
        address,
        nblocks,
        maxtries,
    });
}

export async function getHealth(): Promise<ApiResponse<any>> {
    return invoke<ApiResponse<any>>('cmd_get_health');
}

export async function getLiveness(): Promise<ApiResponse<any>> {
    return invoke<ApiResponse<any>>('cmd_get_liveness');
}

export async function getReadiness(): Promise<ApiResponse<any>> {
    return invoke<ApiResponse<any>>('cmd_get_readiness');
}

export async function getMempool(): Promise<ApiResponse<any>> {
    return invoke<ApiResponse<any>>('cmd_get_mempool');
}

export async function getMempoolTransaction(txid: string): Promise<ApiResponse<any>> {
    return invoke<ApiResponse<any>>('cmd_get_mempool_transaction', { txid });
}

export async function getTransactions(): Promise<ApiResponse<any>> {
    return invoke<ApiResponse<any>>('cmd_get_transactions');
}

export async function getAddressTransactions(address: string): Promise<ApiResponse<any>> {
    return invoke<ApiResponse<any>>('cmd_get_address_transactions', { address });
}

export async function createWalletAdmin(label?: string): Promise<ApiResponse<any>> {
    return invoke<ApiResponse<any>>('cmd_create_wallet_admin', { label });
}

export async function getAddressesAdmin(): Promise<ApiResponse<any>> {
    return invoke<ApiResponse<any>>('cmd_get_addresses_admin');
}

export async function getWalletInfoAdmin(address: string): Promise<ApiResponse<any>> {
    return invoke<ApiResponse<any>>('cmd_get_wallet_info_admin', { address });
}

export async function getBalanceAdmin(address: string): Promise<ApiResponse<any>> {
    return invoke<ApiResponse<any>>('cmd_get_balance_admin', { address });
}

export async function sendTransactionAdmin(
    from: string,
    to: string,
    amount: number,
): Promise<ApiResponse<any>> {
    return invoke<ApiResponse<any>>('cmd_send_transaction_admin', { from, to, amount });
}

export async function setBaseUrl(url: string): Promise<void> {
    return invoke<void>('cmd_set_base_url', { url });
}

export async function setApiKey(key: string): Promise<void> {
    return invoke<void>('cmd_set_api_key', { key });
}

export async function getConfig(): Promise<{ baseUrl: string; apiKey: string }> {
    return invoke<{ baseUrl: string; apiKey: string }>('cmd_get_config');
}
\end{minted}
\end{listing}

\textbf{Walkthrough: \code{invoke()} is the IPC bridge:}

Each TypeScript function calls \code{invoke(command\_name, payload?)} which:

\begin{enumerate}
  \item Serializes the arguments to JSON.
  \item Sends them across the IPC boundary to the Rust backend.
  \item Waits for the command handler to finish.
  \item Deserializes and returns the result.
\end{enumerate}

The pattern is clean: no manual message queues or event systems---just async/await.

% ──────────────────────────────────────────────────────────────────────────────
\section{React Frontend Infrastructure}
\label{sec:ch17-frontend-infra}

The frontend uses React Query (v5) for async state management:

\begin{listing}[htbp]
\caption{React Query hooks (\code{src/hooks/useBlockchain.ts})}
\label{lst:ch17-hooks}
\begin{minted}[fontsize=\footnotesize]{typescript}
import { useQuery } from '@tanstack/react-query';
import * as commands from '../commands';

export function useBlockchainInfo() {
    return useQuery({
        queryKey: ['blockchainInfo'],
        queryFn: () => commands.getBlockchainInfo(),
        staleTime: 10000, // 10 seconds
    });
}

export function useLatestBlocks() {
    return useQuery({
        queryKey: ['latestBlocks'],
        queryFn: () => commands.getLatestBlocks(),
        staleTime: 10000,
    });
}

export function useMiningInfo() {
    return useQuery({
        queryKey: ['miningInfo'],
        queryFn: () => commands.getMiningInfo(),
        staleTime: 30000, // 30 seconds
    });
}
\end{minted}
\end{listing}
